\documentclass[landscape,titlepage]{seminar}
\usepackage{semrot}
\usepackage{fancybox}
\usepackage{fancyhdr}
\usepackage{graphicx}


\title{\shadowbox{Presentation using \LaTeX}} % the document title
%\title{\shadowbox{Presentation using \LaTeX\, seminar class}} % the document title

% the author
\author{Latex Smart \\
         6356 Agricultural Road\\
         University of British Columbia\\
         Vancouver BC \\
         V6T 1Z2}

% the date
\date{\today}



\begin{document}

\begin{slide}
\maketitle
\end{slide}

\begin{slide}
\begin{center}
{\large Outline}
\end{center}
\vskip 0.5cm
\begin{itemize}
 \item The {\tt seminar} document class and {\tt fancybox} and {\tt fancyhdr} package 
 \item {\tt Pdflatex} command and {\tt geometry, hyperref} and {\tt color} package
 \item The {\tt foils} document class with the {\tt FoilTeX} package
 \item Powerpoint-like flashy effects by using {\tt ppower4}
\end{itemize}

\end{slide}

\begin{slide}

\newslide
To create a simple presentation document using \LaTeX, we
can use the {\tt seminar} class. 
The syntax is as follows:

{\small
\begin{verbatim}
\documentclass[landscape]{seminar}
...
\begin{document}
\begin{slide}
...
\newslide
...
\end{slide}
\begin{slide}
...
\newslide
...
\end{slide}
\end{document}
\end{verbatim}
}
\end{slide}


% a new slide environment
\begin{slide}

  \begin{center}
    {\bf Basics of Survival Analysis}
  \end{center}
  Notations:
  \begin{itemize}
    \item{ $T$ = a random variable representing time to event
           (e.g. time to death), $T > 0$ }
    \item{ pdf and cdf: $f(t)$ and $F(t)$ }
    \item{ survivor function: $s(t)$ = $1-F(t)$ }
    \item{ hazard function: $h(t)$,
           which is the probability that an individual dies at
           time t, given that he has survived to time t.}
    \item{ relationship between $h(t)$, $f(t)$, $s(t)$:\\
           $h(t)$ = $\frac{f(t)}{s(t)}$ \\
           $s(t)$ = exp$\{- \int^t_0 h(u) du \}$ \\
           So $f(t)$ = $h(t)$exp$\{- \int^t_0 h(u) du \}$
           }
  \end{itemize}
\end{slide}



\begin{slide}
{\large Include Graphics in Slides}
\begin{center}
\fbox{\includegraphics[height=0.8\textheight]{mypic}}
\end{center}
\newslide

The code for the last slide is:

\begin{verbatim}
\begin{slide}
{\large Include Graphics in Slides}
\begin{center}
\fbox{\includegraphics[height=0.8\textheight]{mypic}}
\end{center}
\newslide
\end{verbatim}

Moreover, we need to add \verb+\usepackage{graphicx}+ in the preamble.


\end{slide}

\begin{slide*}
\begin{center}
{\large Portrait Slides}
\end{center}

The syntax is

\begin{verbatim}
\begin{slide*}
 :
\newslide
 :
\end{slide*}
\end{verbatim}

If you have both landscape and portrait slides in your file, there are two ways to print the
slides:
\begin{itemize}
  \item If your dvi driver supports rotation, then you can print all the slides at once.
For example, by including the package {\tt semrot} in the preamble, you can directly print all
the slides at once.
\newslide
  \item You can first print your landscape slides by putting the command
\begin{verbatim}
  \landscapeonly
\end{verbatim}
in the preamble, and then print the portrait slides by inserting instead the command
\begin{verbatim}
  \portraitonly
\end{verbatim}
and including the {\tt portrait} style option.

\end{itemize}

\end{slide*}

\begin{slide}{{\large \verb|\usepackage{fancybox}|}}

\vskip 0.5cm
\verb|\fbox| \fbox{fbox}
\vskip 0.5cm
\verb|\shadowbox| \shadowbox{shadowbox}

\vskip 0.5cm
\verb|\doublebox| \doublebox{doublebox}

\vskip 0.5cm
\verb|\ovalbox| \ovalbox{ovalbox}

\vskip 0.5cm
\verb|\Ovalbox| \Ovalbox{Ovalbox}

\end{slide}

\begin{slide}{{\large \verb+\usepackage{fancyhdr}+}}
\pagestyle{fancy}
\fancyhf{}
\lhead{left header}
\chead{center header $\alpha$}
\rhead{right header}
\renewcommand{\headrulewidth}{0.5pt}
\lfoot{left footer}
\cfoot{center footer \thepage}
\rfoot{right footer}
\renewcommand{\footrulewidth}{0pt}

\begin{itemize}
  \item \verb+\pagestyle{fancy}+
  \item \verb+\fancyhf{}+ \% delete current setting for header and footer
  \item \verb+\lhead{left header}+
  \item \verb+\chead{center header $\alpha$}+
  \item \verb+\rhead{right header}+
  \item \verb+\renewcommand{\headrulewidth}{0.5pt}+
  \item \verb+\lfoot{left footer}+
  \item \verb+\cfoot{center footer \thepage}+
  \item \verb+\rfoot{right footer}+
  \item \verb+\renewcommand{\footrulewidth}{0pt}+
\end{itemize}

\newslide
\pagestyle{fancy}
\fancyhf{}
\fancyhead[L]{left header}
\fancyhead[C]{center header $\alpha$}
\fancyhead[R]{right header}
\renewcommand{\headrulewidth}{0.5pt}
\fancyfoot[L]{left footer}
\fancyfoot[C]{center footer \thepage}
\fancyfoot[R]{right footer}
\renewcommand{\footrulewidth}{0pt}

Alternatively, we can use
\begin{itemize}
  \item \verb+\pagestyle{fancy}+
  \item \verb+\fancyhf{}+ \% delete current setting for header and footer
  \item \verb+\fancyhead[L]{left header}+
  \item \verb+\fancyhead[C]{center header $\alpha$}+
  \item \verb+\fancyhead[C]{right header}+
  \item \verb+\renewcommand{\headrulewidth}{0.5pt}+
  \item \verb+\fancyfoot[L]{left footer}+
  \item \verb+\fancyfoot[C]{center footer \thepage}+
  \item \verb+\fancyfoot[R]{right footer}+
  \item \verb+\renewcommand{\footrulewidth}{0pt}+
\end{itemize}

\newslide
\fancyhf{}
\renewcommand{\headrulewidth}{0pt}
\renewcommand{\footrulewidth}{0pt}
For double-sided documents, we can use combined options.

For example, \verb+\fancyhead[RE, LO]{\thepage}+ means to place the page number
in the right field of the header on even (left) pages and in the right field of the header on
odd (right) pages.

\begin{tabular}{lll}
\hline
\hline
E & & Even page\\
O & & Odd page\\
L & & Left\\
C & & Center\\
R & & Right\\
H & & Header\\
F & & Footer\\
\hline
\hline
\end{tabular}

\end{slide}

\begin{slide*}
\begin{itemize}
  \item We use {\small \verb+\documentclass[landscape,titlepage]{seminar}+}
  \item syntax
\begin{verbatim}
  \begin{slide} 
     :
    \newslide
     :
  \end{slide} % for landscape slide
  \begin{slide*} 
     :
    \newslide 
     :
  \end{slide*} % for portrait slide
\end{verbatim}
  \item package we use:
    \begin{itemize}
      \item \verb+\usepackage{semrot}+ % rotate slides 
      \item \verb+\usepackage{fancybox}+  
      \item \verb+\usepackage{fancyhdr}+ % fancy header and footer
      \item \verb+\usepackage{graphicx}+
    \end{itemize}
\end{itemize}
\end{slide*}

\begin{slide}{Class and Package}

The difference between {\tt class} and {package} is that
\begin{itemize}
  \item Commands in a class usually cannot be used in other class. For example, \verb+\newslide+
command can only be used in {\tt seminar} class.
  \item Commands in a package can be used in any class.
\end{itemize}

\end{slide}

\end{document}
